\section{Experimental Setup}\label{experimental-setup}

\subsection{Benchmark and Task Suite}\label{benchmark-and-task-suite}

We evaluate on the 82-task OrchestraBench suite from our prior topology
benchmarking study \citep{anonymous_2026}. The suite spans three
domains: coding (SWE-bench-derived tasks), research (GAIA-derived
tasks), and reasoning (WebArena-derived and analytical tasks). Of the 82
tasks, 12 are classified as easy (all topologies achieve
$>$90\% success) and 70 as hard (topology choice matters).
Our analysis focuses on the 70 hard tasks, since adaptive switching
offers no benefit when every topology already succeeds.

Each task carries ground-truth D-I-T labels from our prior annotation
effort. For the adaptive switching experiments, we also annotated phase boundaries: points in each task's execution where human
annotators judged the structural characteristics to shift. Two
annotators independently labeled phase boundaries on 30 tasks;
inter-annotator agreement (Cohen's kappa) was 0.72, indicating
substantial agreement. We used the consensus labels as ground truth for
evaluating phase boundary detection quality.

Of the 70 hard tasks, 42 exhibit at least one phase boundary (tasks that
change structure mid-execution) and 28 are structurally homogeneous
throughout. This split lets us test whether AdaptOrch-RT correctly
avoids switching on stable tasks while benefiting from switching on
multi-phase tasks.

\subsection{Baselines}\label{baselines}

We compare six systems:

\begin{itemize}
\item \textbf{Static-flat}: flat conversation on all tasks.
\item \textbf{Static-hierarchical}: hierarchical decomposition on all tasks.
\item \textbf{Static-debate}: multi-agent debate on all tasks. This was the best single topology in our prior work.
\item \textbf{Static-oracle}: the D-I-T routing classifier from our prior work, which selects the single best topology per task before execution but cannot switch.
\item \textbf{Random switching}: switches topology at random time points with the same average switching frequency as AdaptOrch-RT, to control for the possibility that any switching helps regardless of timing.
\item \textbf{AdaptOrch-RT}: our system.
\end{itemize}

\subsection{Metrics}\label{metrics}

We report four metrics:

\begin{itemize}
\item \textbf{Task success rate}: binary, judged against gold-standard task outputs.
\item \textbf{Switching frequency}: mean number of topology switches per task.
\item \textbf{Switching precision}: fraction of switches that improved the subsequent outcome (the task segment after the switch succeeded when a counterfactual continuation with the prior topology would not have).
\item \textbf{Token overhead}: total tokens consumed, including switching overhead, normalized against the static-oracle baseline.
\end{itemize}

\subsection{Implementation Details}\label{implementation-details}

All experiments use Claude Opus 4.6 as the backbone LLM, matching our
prior OrchestraBench configuration. Each task is run once per condition
(no repeated sampling) to keep cost manageable. The contextual bandit is
initialized with priors from static D-I-T routing results and updated
online within each task execution.

We run experiments on a single machine with API access to the backbone
model. Total experimental cost across all conditions was approximately
4.2M tokens. The sliding window size is \(w = 10\), phase boundary
threshold \(\delta = 0.4\), smoothing parameter \(\alpha = 0.3\),
exploration rate \(\epsilon = 0.1\), and switching gain threshold
\(\gamma = 0.15\). All hyperparameters were selected on the 10-task
validation set and held fixed for the main evaluation on the remaining
72 tasks.
